\subsection{主部与渐近展开}

设 E 是一个由取值于非离散赋值域 K 中的函数组成的比较尺度。设 V 是 K 上的一个赋范空间，f 是 \(\mathcal{H}(\mathfrak{F}, V)\) 中的一个函数；如果存在一个函数 \(g \in E\) 和 V 中的一个元素 \(a \neq 0\) 使得 \(f \sim ag\) ，就说 ag 是 f 相对于尺度 E 的一个主部。由 V, p.~10 的定义 1，f 相对于 E 只能有一个主部，因为若 \(g_{1}\) 、 \(g_{2}\) 是 E 中的两个函数且 \(a_{1}\) 、 \(a_{2}\) 是 V 中两个 \(\neq 0\) 的元素，则关系 \(a_{1}g_{1} \sim a_{2}g_{2}\) 蕴含 \(|g_{1}| \asymp |g_{2}|\) ，从而 \(g_{1} = g_{2}\) ，于是 \((\mathbf{a}_{2} - \mathbf{a}_{1})g_{1} \ll g_{1}\) ，而由于 \(g_{1}\) 在 E 的任何集上都不恒为零，这蕴含 \(a_{2} = a_{1}\) 。

若 f 相对于一个比较尺度 E 有主部，则它相对于每个比较尺度 \(E' \supset E\) 有同一主部。

例. 1) 对实的（相应地，复的）且趋于 \(+\infty\) （相应地， \(\infty\) ）的 x，每个多项式 \(a_{0}x^{n} + a_{1}x^{n-1} + \ldots + a_{n}\) （系数在 V 中，且 \(a_{0} \neq 0\) ）以 \(a_{0}x^{n}\) 为主部（相对于尺度 \(x^{n}\) ，或任何包含 \(x^{n}\) 的尺度）。由此推出，每个有理分式 \(\frac{a_{0}x^{n} + \cdots + a_{m}}{b_{0}x^{n} + \cdots + b_{n}}\) （实或复系数，且 \(a_{0}b_{0} \neq 0\) ）相对于同一尺度以 \(\frac{a_{0}}{b_{0}}x^{m-n}\) 为主部。

\begin{enumerate}
\def\labelenumi{\arabic{enumi})}
\setcounter{enumi}{1}
\tightlist
\item
  一个函数可以与一个尺度的所有函数可比较，却相对于该尺度没有主部。例如，当实数 \(x\) 趋于 \(+\infty\) 时， \(\sqrt{x}\) 相对于尺度 \(x^n\) （其中 \(n\) 为有理整数）没有主部； \(\log x\) 相对于 \(x^\alpha\) （ \(\alpha\) 任意实）的尺度没有主部； \(\exp(\sqrt{\log x})\) 与 \(x^x = e^{x\log x}\) 相对于 \(x^\alpha (\log x)^\beta\) 的尺度没有主部，也不相对于 \(\exp(p(x))\) （ \(p\) 为无常数项多项式）的尺度有主部。
\end{enumerate}

主部的概念容许广泛的推广。设函数 \(\mathbf{f} \in \mathcal{H}(\mathfrak{F}, V)\) 有主部 \(\mathbf{a}_1g_1\) （相对于尺度 \(\mathcal{E}\) ）；关系 \(\mathbf{f} \sim \mathbf{a}_1g_1\) 等价于 \(\mathbf{f} - \mathbf{a}_1g_1 \ll g_1\) (V, p.~216, def. 4)；为更细致地研究函数 \(\mathbf{f}\) ，我们因而被引导去考虑函数 \(\mathbf{f} - \mathbf{a}_1g_1\) 。若这个函数有主部 \(\mathbf{a}_2g_2\) （相对于 \(\mathcal{E}\) ），则必有 \(g_2 \ll g_1\) 且 \(\mathbf{f} - \mathbf{a}_1g_1 - \mathbf{a}_2g_2 \ll g_2\) 。

更一般地，设尺度 E 以参数形式写成 \((g_{\alpha})\) ，其中 \(\alpha\) 取遍一个指标集 A，它被赋予一个同构于 E 的序结构的反序的全序结构：于是关系 \(\alpha < \beta\) 等价于 \(g_{\beta} \ll g_{\alpha}\) 。在这些条件下：

定义 2. 我们说函数 \(\mathbf{f} \in \mathcal{H}(\mathfrak{F}, V)\) 容许一个精确到 \(g_{\alpha}\) 的渐近展开（相对于尺度 \(\mathcal{E}\) ），如果存在 V 的元素的一个族 \((\mathbf{a}_{\lambda})_{\lambda \leqslant \alpha}\) （除有限多个外全为 0），使得 \(\mathbf{f} - \sum_{\lambda \leqslant \alpha} \mathbf{a}_{\lambda} g_{\lambda} \ll g_{\alpha}\) 。我们说 \(\sum_{\lambda \leqslant \alpha} \mathbf{a}_{\lambda} g_{\lambda}\) 是 \(\mathbf{f}\) 的一个精确到 \(g_{\alpha}\) 的渐近展开， \(\mathbf{a}_{\lambda} g_{\lambda} (\lambda \leqslant \alpha)\) 是它的项， \(\mathbf{a}_{\lambda}\) 是它的系数，而函数 \(\mathbf{r}_{\alpha} = \mathbf{f} - \sum_{\lambda \leqslant \alpha} \mathbf{a}_{\lambda} g_{\lambda}\) 是这个展开的余项。

为表示 \(\sum_{\lambda \leqslant \alpha} \mathbf{a}_{\lambda} g_{\lambda}\) 是 \(\mathbf{f}\) 的一个精确到 \(g_{\alpha}\) 的渐近展开这一事实，最常见的是只写

\[
\mathbf {f} = \sum_ {\lambda \leqslant \alpha} \mathbf {a} _ {\lambda} g _ {\lambda} + o (g _ {\alpha}) \quad (\text { or } \mathbf {f} = \sum_ {\lambda \leqslant \alpha} \mathbf {a} _ {\lambda} g _ {\lambda} + o _ {k} (g _ {\alpha})
\]

如果证明中有多个函数），遵循 V, p.~219 与 220 的记号。

对于相对于同一尺度 E 的两个渐近展开（两个函数的，相同与否），我们说精确度指标较大的那个更精确。

若 \(\sum_{\lambda \leqslant \alpha} \mathbf{a}_{\lambda} g_{\lambda}\) 是 \(\mathbf{f}\) 的一个精确到 \(g_{\alpha}\) 的渐近展开，则对每个 \(\beta < \alpha\) ， \(\sum_{\lambda \leqslant \beta} \mathbf{a}_{\lambda} g_{\lambda}\) 是 \(\mathbf{f}\) 的一个精确到 \(g_{\beta}\) 的渐近展开 (V, p.~215, prop. 5)：说它是通过把给定展开 \(\sum_{\lambda \leqslant \alpha} \mathbf{a}_{\lambda} g_{\lambda}\) （ \(\mathbf{f}\) 的）的精确度降到 \(g_{\beta}\) 而得到的。

若 \(\sum_{\lambda \leqslant \alpha} \mathbf{a}_{\lambda} g_{\lambda}\) 与 \(\sum_{\lambda \leqslant \alpha} \mathbf{b}_{\lambda} g_{\lambda}\) 是同一精确度 \(g_{\alpha}\) 的渐近展开（两个函数 \(\mathbf{f}_1, \mathbf{f}_2\) 的），则 \(\sum_{\lambda \leqslant \alpha} (\mathbf{a}_{\lambda} + \mathbf{b}_{\lambda}) g_{\lambda}\) 是 \(\mathbf{f}_1 + \mathbf{f}_2\) 的一个精确到

精确度 \(g_{\alpha}\) (V, p.~215, prop. 5)；且对每个标量 \(c\) ， \(\sum_{\lambda \leqslant \alpha} \mathbf{a}_{\lambda}cg_{\lambda}\) 是 \(\mathbf{f}_1c\) 的一个精确到 \(g_{\alpha}\) 的渐近展开。由此推出，若一个函数容许一个精确到 \(g_{\alpha}\) 的渐近展开，则这个展开是唯一的：只需看到函数 0 不容许一个精确度 \(g_{\alpha}\) 的、具有 \(\neq 0\) 系数的渐近展开。因为，若 \(0 = \sum_{\lambda \leqslant \alpha} \mathbf{a}_{\lambda}g_{\lambda} + \mathbf{r}_{\alpha}\) ，且 \(\gamma\) 是指标 \(\lambda \leqslant \alpha\) （使 \(\mathbf{a}_{\lambda} \neq 0\) 者）中最小者，则将有 \(\mathbf{a}_{\gamma}g_{\gamma} = -\sum_{\gamma < \lambda \leqslant \alpha} \mathbf{a}_{\lambda}g_{\lambda} - \mathbf{r}_{\alpha} \ll g_{\gamma}\) ，这是荒谬的。

说一个函数 \(\mathbf{f}\) 容许一个精确到 \(g_{\alpha}\) 的渐近展开且其系数全为零，等价于说 \(\mathbf{f} \ll g_{\alpha}\) 。若 \(\mathbf{f}\) 容许一个渐近展开 \(\sum_{\lambda \leqslant \alpha} \mathbf{a}_{\lambda} g_{\lambda}\) （精确到 \(g_{\alpha}\) ），其系数不全为零，且 \(\gamma\) 是指标 \(\lambda\) （使 \(\mathbf{a}_{\lambda} \neq 0\) 者）中最小者，则 \(\mathbf{a}_{\gamma} g_{\gamma}\) 是 \(\mathbf{f}\) 相对于尺度 \(\mathcal{E}\) 的主部，因为有 \(\mathbf{f} - \mathbf{a}_{\gamma} g_{\gamma} = \sum_{\gamma < \lambda \leqslant \alpha} \mathbf{a}_{\lambda} g_{\lambda} + \mathbf{r}_{\alpha} \ll g_{\gamma}\) ；类似地，若 \(\mu \leqslant \alpha\) 是使 \(\mathbf{a}_{\mu} \neq 0\) 的一个指标，则 \(\mathbf{a}_{\mu} g_{\mu}\) 是 \(\mathbf{f} - \sum_{\lambda < \mu} \mathbf{a}_{\lambda} g_{\lambda}\) 的主部。

在应用中最重要的渐近展开是相对于 \(x^{-n}\) （相应地， \(z^{-n}\) ）的尺度的那些，其中 n 是正或负整数，当 x 趋于 \(+\infty\) 或 \(-\infty\) （相应地，当复数 z 趋于 \(\infty\) ）时；或相对于 \((x - c)^{n}\) （相应地， \((z - c)^{n}\) ）的尺度的那些，当实数 x 从右边或左边趋于 c（相应地，当复数 z 趋于 c）时。我们在 I, p.~21 中已看到，在一点 \(c \in R\) 处 k 次可导的每个实变量 x 的向量函数在该点容许一个 k 阶泰勒展开，即一个精确到 \((x - c)^{k}\) 的、相对于 \((x - c)^{n}\) 的尺度的渐近展开。

